\documentclass[11pt,a4paper]{article}

% --- Variables du TP
\newcommand{\TPModule}{Datawarehouse}
\newcommand{\TPNumero}{2}
\newcommand{\TPKeywords}{datawarehouse, datamart, star schema, dbt, KPI}
\newcommand{\TPSujet}{Modélisation décisionnelle et couche GOLD}
\newcommand{\TPTheme}{Clinique}
\usepackage[utf8]{inputenc}
\usepackage[french]{babel}
\usepackage[T1]{fontenc}
\usepackage{amsmath}
\usepackage{amsfonts}
\usepackage{xcolor}
\usepackage{amssymb}
\usepackage{graphicx}
\usepackage[left=2cm,right=2cm,top=2cm,bottom=2cm]{geometry}
\usepackage{fancyhdr}
\usepackage[bottom]{footmisc}
\usepackage{tikz}
\usepackage{fourier}
\usepackage{pst-text}
\usepackage{multicol}
\usepackage{comment}
\usepackage{array}
\usepackage{enumitem}
\setlist[itemize]{label=\textbullet} % pour ne plus utiliser les bulletpoints après \item
\setlength{\parindent}{0pt} % pour supprimer les \noindent
\usepackage{amssymb} % pour les symboles
\usepackage{tabularx}
\usepackage{tcolorbox}
\usepackage{fontawesome}
\usepackage{listings}

% --- Réglages listings (inchangé)
\lstset{
  basicstyle=\ttfamily\small,
  backgroundcolor=\color{gray!10},
  frame=single,
  rulecolor=\color{gray!40},
  breaklines=true,
  columns=fullflexible,
  showstringspaces=false,
  aboveskip=0.5em,
  belowskip=0.5em,
  keywordstyle=\color{blue},
  commentstyle=\color{gray},
}
\lstdefinelanguage{markdown}{
  sensitive=true,
  morecomment=[l]{<!--},
  morecomment=[s]{<!--}{-->},
  morekeywords={},
}

% --- Réglages listings (inchangé)
\lstset{
  basicstyle=\ttfamily\small,
  backgroundcolor=\color{gray!10},
  frame=single,
  rulecolor=\color{gray!40},
  breaklines=true,
  columns=fullflexible,
  showstringspaces=false,
  aboveskip=0.5em,
  belowskip=0.5em,
  keywordstyle=\color{blue},
  commentstyle=\color{gray},
}
\lstdefinelanguage{markdown}{
  sensitive=true,
  morecomment=[l]{<!--},
  morecomment=[s]{<!--}{-->},
  morekeywords={},
}

% --- Fancyhdr construit depuis les variables
\pagestyle{fancy}
\setlength{\headheight}{26pt}
\renewcommand{\baselinestretch}{1.3}

\fancyhead[L]{\textbf{EPSI\\\TPModule}}
\fancyhead[R]{\textbf{TP n°\TPNumero\\\TPSujet}}
\renewcommand{\arraystretch}{1}

% --- Hyperref + metadata construits depuis les variables
\usepackage[
  colorlinks=true,
  linkcolor=blue!70!black,
  urlcolor=blue!70!black,
  citecolor=blue!70!black
]{hyperref}

\hypersetup{
  pdftitle={TP n°\TPNumero\space - \TPSujet},
  pdfauthor={Thibault CLEMENT},
  pdfsubject={TP n°\TPNumero\space - \TPSujet},
  pdfkeywords={\TPKeywords},
}

% --- Composants réutilisables
\newcommand{\TPFrontPage}{%
  ~
  \begin{center}
    \framebox{\parbox[c]{11cm}{%
      \begin{center}
        \Large{\textbf{TP n°\TPNumero~:~\TPSujet\\\large{\TPTheme}}}
      \end{center}
    }}
  \end{center}
  \vspace{0.5cm}
}

% --- Tes macros custom existantes
\newcommand{\file}[1]{~\texttt{\textcolor{blue!70!black}{#1}}~}

\newtcolorbox{reflectionbox}[1][]{%
  colback=blue!5,
  colframe=blue!60,
  title={\faLightbulbO~~Réflexions},
  fonttitle=\bfseries,
  left=1mm,
  right=1mm,
  top=1mm,
  bottom=1mm,
  #1
}


\newcommand{\checkbox}{\(\square\)~~}
\graphicspath{{images/}}
\setlength{\columnseprule}{1pt}
\def\columnseprulecolor{\color{black}}



\begin{document}
\TPFrontPage

% ---------------------------------------------------------------------------------------------

\section*{Contexte}

Lors du TP1, vous avez :

\begin{itemize}
  \item exploré les données sources,
  \item mis en place une couche \textbf{BRONZE} historisée,
  \item construit une couche \textbf{SILVER} propre et testée via dbt.
\end{itemize}

Vous devez maintenant concevoir la \textbf{couche GOLD},
c’est-à-dire un \textbf{datamart décisionnel} permettant de répondre
aux questions stratégiques du directeur.

\bigskip

\section*{Rappel : Questions métier du directeur}

L’objectif final est de pouvoir répondre aux questions suivantes :

\begin{enumerate}
  \item Quels services sont le plus souvent en tension ?
  \item Quel est le taux moyen d’occupation par service ?
  \item Combien de patients sont refusés chaque semaine et pourquoi ?
  \item La satisfaction diminue-t-elle en période de forte activité ?
  \item Existe-t-il un lien entre présence du personnel et tension ?
  \item Quelle est la durée moyenne de séjour par service ?
  \item Quels services semblent les plus rentables ?
  \item Peut-on anticiper les semaines critiques ?
\end{enumerate}

\textbf{Objectif du TP2 :}
construire le modèle décisionnel permettant de produire ces indicateurs.

% ---------------------------------------------------------------------------------------------

\section*{Architecture technique}

Base de données : \file{dwh/clinic\_dwh.db}

Séparation des couches par préfixe :

\begin{itemize}
  \item \file{bronze\_*}
  \item \file{silver\_*}
  \item \file{gold\_*} (à construire)
\end{itemize}

Les modèles GOLD devront être implémentés via \textbf{dbt}
dans un dossier \file{models/gold/}.

% ---------------------------------------------------------------------------------------------

\section{Étape 1 : Définition du fait principal et du grain}

\subsection*{Objectif}

Définir précisément ce que l’on souhaite analyser.

\subsection*{Travail demandé}

\begin{enumerate}
  \item Identifier le \textbf{fait métier principal}.
  \item Définir son \textbf{grain exact}.
  \item Justifier ce choix.
  \item Lister les mesures associées.
\end{enumerate}

\textbf{Indication :}
un grain pertinent est souvent :
\begin{center}
\textit{1 ligne = 1 service × 1 semaine}
\end{center}

% ---------------------------------------------------------------------------------------------

\section{Étape 2 : Construction des indicateurs hebdomadaires}

Aucune table agrégée hebdomadaire n’est fournie.
Vous devez la construire à partir de la couche \textbf{silver}.

\subsection*{Objectif}

Créer une table décisionnelle agrégée par \textbf{service} et \textbf{yearweek}.

\subsection*{Indicateurs attendus (minimum)}

\begin{itemize}
  \item Nombre de patients présents.
  \item Nombre de demandes.
  \item Nombre de patients admis.
  \item Nombre de patients refusés.
  \item Taux de refus.
  \item Durée moyenne de séjour.
  \item Satisfaction moyenne.
  \item Nombre de membres du personnel présents.
\end{itemize}

\subsection*{Travail demandé}

\begin{enumerate}
  \item Identifier les tables silver nécessaires.
  \item Gérer les jointures et éventuelles duplications.
  \item Créer un modèle dbt : \file{gold\_fact\_service\_weekly}.
  \item Documenter les calculs effectués.
\end{enumerate}

\subsection*{Points d’attention}

\begin{itemize}
  \item Attention au grain lors des jointures.
  \item Vérifier la cohérence : 
  \file{patients\_refused = requests - admitted}.
  \item Ne pas utiliser d’information future pour calculer une métrique passée.
\end{itemize}

% ---------------------------------------------------------------------------------------------

\section{Étape 3 : Construction des dimensions}

\subsection*{Objectif}

Mettre en place un schéma en étoile.

\subsection*{Dimensions minimales attendues}

\begin{itemize}
  \item \file{gold\_dim\_service}
  \item \file{gold\_dim\_date} (ou week)
  \item \file{gold\_dim\_staff} (si pertinent)
\end{itemize}

\subsection*{Travail demandé}

\begin{enumerate}
  \item Définir les attributs de chaque dimension.
  \item Identifier les clés primaires.
  \item Relier les dimensions à la table de faits.
  \item Justifier les cardinalités.
\end{enumerate}

% ---------------------------------------------------------------------------------------------

\section{Étape 4 : Implémentation GOLD avec dbt}

\subsection*{Consignes}

\begin{enumerate}
  \item Implémenter les modèles GOLD dans \file{models/gold/}.
  \item Matérialisation en tables.
  \item Documenter les modèles.
\end{enumerate}

\subsection*{Tests obligatoires}

\begin{itemize}
  \item \file{unique} + \file{not\_null} sur les PK.
  \item \file{relationships} entre faits et dimensions.
  \item Test métier simple sur un KPI.
\end{itemize}

\subsection*{Commandes attendues}

\begin{lstlisting}
dbt run  --select gold
dbt test --select gold
\end{lstlisting}

% ---------------------------------------------------------------------------------------------

\section{Étape 5 : Historisation et volumétrie}

\subsection*{Réflexion}

\begin{enumerate}
  \item Estimer la volumétrie annuelle.
  \item Identifier une dimension susceptible d’évoluer.
  \item Proposer une stratégie SCD (Type 1 ou Type 2). Justifier le choix.
\end{enumerate}

% ---------------------------------------------------------------------------------------------

\section{Étape 6 : Exploitation décisionnelle (Dashboard BI)}

\subsection*{Objectif}

Valider la cohérence du datamart GOLD en construisant un premier
tableau de bord décisionnel permettant de répondre aux questions
du directeur.

\subsection*{Outil recommandé}

Power BI (ou outil équivalent : Tableau, Superset, etc.).

\subsection*{Travail demandé}

\begin{enumerate}
  \item Se connecter à la base \file{clinic\_dwh.db}.
  \item Importer uniquement les tables \file{gold\_*}.
  \item Vérifier les relations entre tables (PK/FK).
\end{enumerate}

\subsection*{KPI minimums attendus}

Construire au minimum les indicateurs suivants :

\begin{itemize}
  \item Taux d’occupation moyen par service.
  \item Nombre de patients refusés par semaine.
  \item Taux de refus.
  \item Satisfaction moyenne.
  \item Durée moyenne de séjour.
  \item Nombre de membres du personnel présents.
\end{itemize}

Pour chaque KPI, préciser :

\begin{itemize}
  \item la définition métier,
  \item la formule utilisée,
  \item le niveau d’agrégation.
\end{itemize}

\subsection*{Structure attendue du dashboard}

Créer au minimum :

\begin{itemize}
  \item Une page \textbf{Pilotage global} :
  \begin{itemize}
    \item KPI principaux (cartes),
    \item évolution temporelle,
    \item filtre par service.
  \end{itemize}
  \item Une page \textbf{Analyse détaillée} :
  \begin{itemize}
    \item comparaison entre services,
    \item focus sur les refus et la satisfaction,
    \item possibilité de filtrer par période.
  \end{itemize}
\end{itemize}

\subsection*{Analyse attendue}

Rédiger une courte analyse (10 lignes environ) répondant à au moins
\textbf{3 questions du directeur} à partir du dashboard construit.

\subsection*{Livrables attendus}

\begin{itemize}
  \item Fichier Power BI (\file{.pbix}) ou captures d’écran.
  \item Liste des KPI définis (avec formules).
  \item Analyse rédigée.
\end{itemize}

\end{document}